\section{Derive the continuum level-density relation}
\label{app:cld-derivation}
% BEGIN CURATED SUBJECT INDEX
\index{resolvent}
\index{trace-class operator}
\index{spectral shift function}
% END CURATED SUBJECT INDEX

For a trace-class resolvent difference, define the perturbation determinant
\begin{equation}
  \mathcal D(z)
  =\det\left[(z-H)(z-H_0)^{-1}\right] .
  \label{eq:perturbation-determinant}
\end{equation}
Its logarithmic derivative is
\begin{equation}
  \frac{\rmd}{\rmd z}\log\mathcal D(z)
  =\Tr\left[(z-H)^{-1}-(z-H_0)^{-1}\right] .
  \label{eq:determinant-resolvent-trace}
\end{equation}
Taking the imaginary part of the upper boundary value gives
Eq.~\eqref{eq:cld-physical}.  The ratio of boundary values of
$\mathcal D$ is related to $\det S$; with
$\det S=\rme^{2\rmi\vsub{\delta}{tot}}$, differentiation gives
Eq.~\eqref{eq:birman-krein-differential}.

Complex deformation does not change $\mathcal D(z)$ in the analytic region
between the original and rotated contours.  This gives
Eq.~\eqref{eq:cld-csm}.  In a finite basis, the determinant is a ratio of
characteristic polynomials, whose logarithmic derivative gives the sum in
Eq.~\eqref{eq:finite-cld}.

Bound states contribute jumps to the integrated spectral shift.  Care is
therefore needed when comparing phases at threshold; Levinson-type constants
depend on half-bound states, dimensions, and long-range tails.
